\begin{codebox}
\codeline{\textcolor{cmtgray}{\#\ Text2SPARQL：自然语言到查询语言的安全转换}}
\codeline{\textcolor{cmtgray}{\#\ Text2SPARQL:\ Safe\ Natural\ Language\ to\ Query\ Translation}}
\codeline{\textcolor{cmtgray}{\#}}
\codeline{\textcolor{cmtgray}{\#\ 让LLM"写查询"而不是"答事实"——事实只能来自查询结果}}
\codeblank{}
\codeline{\textcolor{cmtgray}{\#\ ==============================}}
\codeline{\textcolor{cmtgray}{\#\ 1.\ 三种技术路线}}
\codeline{\textcolor{cmtgray}{\#\ ==============================}}
\codeline{\textcolor{cmtgray}{\#\ 路线A\ 模板槽填充：}}
\codeline{\textcolor{cmtgray}{\#\ \ \ 预定义查询模板，LLM只负责识别意图与填槽}}
\codeline{\textcolor{cmtgray}{\#\ \ \ 高可控、零注入风险；覆盖面有限}}
\codeline{\textcolor{cmtgray}{\#\ 路线B\ LLM直接生成SPARQL\ +\ 校验层：}}
\codeline{\textcolor{cmtgray}{\#\ \ \ 灵活覆盖长尾问题；必须配校验层（见第3节）}}
\codeline{\textcolor{cmtgray}{\#\ 路线C\ 语义解析（训练专用模型）：}}
\codeline{\textcolor{cmtgray}{\#\ \ \ 精度高；需要标注数据，迭代成本高}}
\codeline{\textcolor{cmtgray}{\#\ 工程组合：高频问题走A，长尾走B，A/B共享同一校验层}}
\codeblank{}
\codeline{\textcolor{cmtgray}{\#\ ==============================}}
\codeline{\textcolor{cmtgray}{\#\ 2.\ 路线A示例：模板槽填充}}
\codeline{\textcolor{cmtgray}{\#\ ==============================}}
\codeline{\textcolor{cmtgray}{\#\ 模板T01\ 设备状态查询：}}
\codeline{\textcolor{cmtgray}{\#\ \ \ 意图关键词：状态/在干嘛/空闲吗}}
\codeline{\textcolor{cmtgray}{\#\ \ \ 模板：}}
\codeline{\textcolor{cmtgray}{\#\ \ \ \ \ SELECT\ ?status\ WHERE\ \{\ \{EQUIPMENT\}\ mfg:hasStatus\ ?status\ \}}}
\codeline{\textcolor{cmtgray}{\#\ \ \ 槽位：\{EQUIPMENT\}\ ←\ 实体链接结果（别名表见第7章）}}
\codeblank{}
\codeline{\textcolor{cmtgray}{\#\ 模板T02\ 能力查询：}}
\codeline{\textcolor{cmtgray}{\#\ \ \ "哪些设备能加工\{MATERIAL\}？"}}
\codeline{\textcolor{cmtgray}{\#\ \ \ \ \ SELECT\ ?e\ WHERE\ \{}}
\codeline{\textcolor{cmtgray}{\#\ \ \ \ \ \ \ \ \ ?e\ rdf:type/rdfs:subClassOf*\ mfg:Equipment\ ;}}
\codeline{\textcolor{cmtgray}{\#\ \ \ \ \ \ \ \ \ \ \ \ mfg:canProcess\ \{MATERIAL\}\ .\ \}}}
\codeblank{}
\codeline{\textcolor{cmtgray}{\#\ 填槽过程：}}
\codeline{\textcolor{cmtgray}{\#\ \ \ "3号车床现在空闲吗"\ \(\rightarrow\)\ 意图T01\ +\ 实体mfg:Lathe\_003}}
\codeline{\textcolor{cmtgray}{\#\ \ \ \(\rightarrow\)\ SELECT\ ?status\ WHERE\ \{\ mfg:Lathe\_003\ mfg:hasStatus\ ?status\ \}}}
\codeline{\textcolor{cmtgray}{\#\ \ \ \(\rightarrow\)\ 返回\ mfg:Idle\ \(\rightarrow\)\ 生成回答"是，当前空闲"}}
\codeblank{}
\codeline{\textcolor{cmtgray}{\#\ ==============================}}
\codeline{\textcolor{cmtgray}{\#\ 3.\ 路线B的校验层（关键防线）}}
\codeline{\textcolor{cmtgray}{\#\ ==============================}}
\codeline{\textcolor{cmtgray}{\#\ LLM生成的SPARQL必须依次通过四道检查：}}
\codeblank{}
\codeline{\textcolor{cmtgray}{\#\ (1)\ 语法校验}}
\codeline{\textcolor{cmtgray}{\#\ \ \ \ \ 用SPARQL解析器（Jena\ ARQ\ /\ rdflib）解析，失败\(\rightarrow\)重试或回退模板}}
\codeblank{}
\codeline{\textcolor{cmtgray}{\#\ (2)\ 操作白名单（防注入/防破坏）}}
\codeline{\textcolor{cmtgray}{\#\ \ \ \ \ 只允许\ SELECT\ /\ ASK；}}
\codeline{\textcolor{cmtgray}{\#\ \ \ \ \ 拒绝\ INSERT\ /\ DELETE\ /\ DROP\ /\ LOAD\ /\ SERVICE（防外联）}}
\codeline{\textcolor{cmtgray}{\#\ \ \ \ \ 生产端点同时在服务端禁用update（双保险）}}
\codeblank{}
\codeline{\textcolor{cmtgray}{\#\ (3)\ Schema校验（防"词汇幻觉"）}}
\codeline{\textcolor{cmtgray}{\#\ \ \ \ \ 查询中出现的类/属性必须存在于本体：}}
\codeline{\textcolor{cmtgray}{\#\ \ \ \ \ \ \ LLM编造\ mfg:hasMaxSpeed（本体中只有\ mfg:hasSpindleSpeed）}}
\codeline{\textcolor{cmtgray}{\#\ \ \ \ \ \ \ \(\rightarrow\)\ Schema校验失败\ \(\rightarrow\)\ 把本体属性清单回喂LLM重新生成}}
\codeline{\textcolor{cmtgray}{\#\ \ \ \ \ 这是Text2SPARQL最常见错误，校验后准确率提升显著}}
\codeblank{}
\codeline{\textcolor{cmtgray}{\#\ (4)\ 资源限额}}
\codeline{\textcolor{cmtgray}{\#\ \ \ \ \ 强制\ LIMIT（如1000）、查询超时（如5s）、禁用笛卡尔积形态}}
\codeblank{}
\codeline{\textcolor{cmtgray}{\#\ ==============================}}
\codeline{\textcolor{cmtgray}{\#\ 4.\ 提示词要点（路线B）}}
\codeline{\textcolor{cmtgray}{\#\ ==============================}}
\codeline{\textcolor{cmtgray}{\#\ 提供给LLM的上下文应包含：}}
\codeline{\textcolor{cmtgray}{\#\ \ \ 本体Schema摘要：类层次\ +\ 属性清单（域/值域）+\ 3\textasciitilde{}5个示例查询}}
\codeline{\textcolor{cmtgray}{\#\ \ \ 明确规则："只用清单中的词汇；不确定就返回\ NEED\_CLARIFY"}}
\codeline{\textcolor{cmtgray}{\#\ 示例few-shot（一问一查询），覆盖：单实体查询/过滤/聚合/路径}}
\codeblank{}
\codeline{\textcolor{cmtgray}{\#\ ==============================}}
\codeline{\textcolor{cmtgray}{\#\ 5.\ 失败回退策略}}
\codeline{\textcolor{cmtgray}{\#\ ==============================}}
\codeline{\textcolor{cmtgray}{\#\ 生成失败/校验不过\ \(\rightarrow\)\ 重试（带错误信息）最多2次}}
\codeline{\textcolor{cmtgray}{\#\ \ \ \(\rightarrow\)\ 仍失败\ \(\rightarrow\)\ 回退模板匹配}}
\codeline{\textcolor{cmtgray}{\#\ \ \ \(\rightarrow\)\ 模板也无\ \(\rightarrow\)\ 澄清式反问用户（"您是想查状态还是维护记录？"）}}
\codeline{\textcolor{cmtgray}{\#\ 禁止：在任何失败分支让LLM"凭印象"直接回答事实}}
\codeblank{}
\codeline{\textcolor{cmtgray}{\#\ ==============================}}
\codeline{\textcolor{cmtgray}{\#\ 6.\ 评测指标}}
\codeline{\textcolor{cmtgray}{\#\ ==============================}}
\codeline{\textcolor{cmtgray}{\#\ 语法正确率：可解析的比例（基线，通常很高）}}
\codeline{\textcolor{cmtgray}{\#\ Schema命中率：词汇全部在本体内的比例（核心质量指标）}}
\codeline{\textcolor{cmtgray}{\#\ 执行准确率：查询结果与标准答案一致的比例（最终指标）}}
\codeline{\textcolor{cmtgray}{\#\ 评测集：把第3章CQ清单转成"自然语言问法\(\rightarrow\)标准查询"对，一举两用}}
\end{codebox}
